\section{Introduction}\label{sec:intro}
Spatial transcriptomics platforms differ in almost every respect: imaging-based assays such as Xenium and CosMx
resolve single cells against targeted panels, while sequencing-based assays such as Visium HD and Stereo-seq capture
larger gene sets on grids. When serial sections of one tissue are profiled on several
platforms~\citep{ren2025spatch}, the same cell types land in different regions of an embedding space. Foundation
models~\citep{cui2024scgpt,wang2025scgptspatial,tejadalapuerta2025nicheformer} give every platform a common
vocabulary, but not a common geometry.

A common workflow annotates one high-quality reference and brings other data sets into it. Most integration
methods cannot support this. Harmony~\citep{korsunsky2019harmony}, Scanorama~\citep{hie2019scanorama} and spatially
aware methods such as GraphST~\citep{long2023graphst} return a new embedding for all inputs. The reference moves, and
every classifier, index or annotation built on it must be rebuilt. USHER~\citep{pratapa2025usher} instead keeps the
target fixed and learns a low-capacity map of the source into it, using fused Gromov--Wasserstein (FGW) transport to
find correspondences. Its structure matrices, however, are recomputed in every EM iteration, and this dominates its
run time. Because composite scores are dominated by batch mixing, which moment matching recovers cheaply, we read
results as a trade-off between mixing and conservation.

Our contributions are: (i) \frost, a post-hoc, reference-preserving FGW method whose structure matrices are built
once per batch ($2B$ rather than $2BT$ computations for $B$ batches and $T$ iterations) and which ends with
closed-form moment matching; (ii) a 12-pair benchmark across three tumour types and four source platforms against
ten baselines, with a measure of reference drift, in which a spatially aware integration method fails; (iii) a comparison with USHER showing similar total scores from opposite sides of the trade-off at
25-fold fewer structure computations at USHER's 25 EM iterations (a $T$-fold reduction for $T$ iterations); and (iv) a mechanistic account of why global alignment blurs gene-level
neighbourhoods, which sets the method's scope.
